\subsubsection{Examples of Banach Spaces}\label{subsec:banach-space-examples}

\Needspace{4\baselineskip}
\begin{Example}
\label{ex:bounded-continuous-functions-banach}
Let $\Omega$ be an open set in $\mathbb{R}^m$ and let $C_b(\overline\Omega)$ be the set of all bounded continuous functions on $\overline\Omega$. We define
\begin{equation}\label{eq:bounded-continuous-functions-supremum-norm}
\|u\|=\sup_{\bm{x}\in\overline\Omega}|u(\bm{x})|.
\end{equation}
Then $C_b(\overline\Omega)$ is a Banach space. Indeed, the properties of a norm are easy to check, and the completeness is a restatement of the well-known fact that uniform limits of continuous functions are continuous (note that convergence in $C_b(\overline\Omega)$ means uniform convergence).
\end{Example}

\Needspace{4\baselineskip}
\begin{Example}
\label{ex:bounded-analytic-functions-banach}
Let $\Omega$ be an open set in $\mathbb{C}$ and let $A(\Omega)$ be the set of all bounded analytic functions on $\Omega$. We define
\begin{equation}\label{eq:bounded-analytic-functions-supremum-norm}
\|u\|=\sup_{z\in\Omega}|u(z)|.
\end{equation}
Then $A(\Omega)$ is a Banach space. Indeed, the fact that uniform limits of analytic functions are analytic is an easy consequence of Morera\textquotesingle s theorem.
\end{Example}

\begin{Example}
\label{ex:continuous-functions-lp-norm}
Let $\Omega$ be an open, locally Jordan measurable (cf. \cite{Bu}) (this means that the intersection with any ball is Jordan measurable) set in $\mathbb{R}^m$ and let $1\leq p<\infty$. Let $\widetilde L^p(\Omega)$ be the set of all continuous functions $u:\overline\Omega\to\mathbb{R}$ for which
\begin{equation}\label{eq:continuous-functions-lp-integral-norm}
\|u\|_p:=\left(\int_\Omega|u(\bm{x})|^p\,\mathrm{d}\bm{x}\right)^{1/p}
\end{equation}
is finite. We want to show that $\widetilde L^p(\Omega)$ is a normed vector space. The first two properties of a norm are trivial, and the triangle inequality is trivial for $p=1$. The triangle inequality for $p>1$ takes some preparation.
\end{Example}

\begin{Lemma}
\label{lem:young-conjugate-exponents-scalar-inequality}
Let $a,b\geq0$, $p,q>1$ and $p^{-1}+q^{-1}=1$. Then
\begin{equation}\label{eq:young-conjugate-exponents-scalar-inequality}
ab\leq\dfrac{a^p}p+\dfrac{b^q}q.
\end{equation}

\end{Lemma}

The proof is obtained by seeking the minimum of the function \[b\mapsto a^p/p+b^q/q-ab.\] See Problem~\ref{ex:young-conjugate-exponents-inequality}. As a consequence of Lemma~\ref{lem:young-conjugate-exponents-scalar-inequality}, we shall now establish Hölder\textquotesingle s inequality.

\begin{Theorem}
\label{thm:continuous-functions-holder-inequality}
Let $f\in\widetilde L^p(\Omega)$, $g\in\widetilde L^q(\Omega)$, where $p,q\in(1,\infty)$ with $p^{-1}+q^{-1}=1$. Then $fg\in\widetilde L^1(\Omega)$, and
\begin{equation}\label{eq:continuous-functions-holder-inequality}
\|fg\|_1\leq\|f\|_p\|g\|_q.
\end{equation}

\end{Theorem}

\begin{Proof}
We only rule out trivial cases by assuming that $\|f\|_p\ne0$, $\|g\|_q\ne0$. By the previous lemma, we find
\begin{equation}\label{eq:holder-normalized-young-pointwise-bound}
\dfrac{|f(\bm{x})|}{\|f\|_p}\dfrac{|g(\bm{x})|}{\|g\|_q}\leq\dfrac1p\dfrac{|f(\bm{x})|^p}{\|f\|_p^p}+\dfrac1q\dfrac{|g(\bm{x})|^q}{\|g\|_q^q}.
\end{equation}
The theorem follows by integrating over $\Omega$.
\end{Proof}

\begin{Theorem}
\label{thm:continuous-functions-lp-triangle-inequality}
Let $f,g\in\widetilde L^p(\Omega)$. Then $f+g\in\widetilde L^p(\Omega)$ and $\|f+g\|_p\leq\|f\|_p+\|g\|_p$.
\end{Theorem}

\begin{Proof}
We note that
\begin{equation}\label{eq:lp-sum-integrability-bound}
|f+g|^p\leq(|f|+|g|)^p\leq2^p(|f|^p+|g|^p);
\end{equation}
hence $f+g\in\widetilde L^p(\Omega)$. Moreover, we have
\begin{equation}\label{eq:lp-triangle-pointwise-product-bound}
|f+g|^p\leq|f+g|^{p-1}|f|+|f+g|^{p-1}|g|,
\end{equation}
and with $p^{-1}+q^{-1}=1$, we have $(p-1)q=p$; hence $|f+g|^{p-1}\in\widetilde L^q(\Omega)$. We can therefore apply Hölder\textquotesingle s inequality and obtain
\begin{equation}\label{eq:lp-triangle-holder-integral-bound}
\int_\Omega|f+g|^p\,\mathrm{d}\bm{x}\leq\left(\int_\Omega|f+g|^p\,\mathrm{d}\bm{x}\right)^{1/q}(\|f\|_p+\|g\|_p).
\end{equation}
From this the triangle inequality follows immediately.
\end{Proof}
